\subsection*{The Croqtile Function}

Start with the signature, which declares more than a C++ function prototype would: every parameter and the result carry a shape.

\begin{lstlisting}[language=croqtile,caption={A kernel signature}]
__co__ s32 [4, 8] ele_add(s32 [4, 8] lhs, s32 [4, 8] rhs) {
\end{lstlisting}

The \texttt{\_\_co\_\_} prefix marks this as a Croqtile function. Unlike regular C++ functions, the signature carries full shape information: \texttt{s32 [4, 8]} means "a 2D tensor of shape 4 \(\times\) 8 with element type \texttt{s32} (signed 32-bit integer)". Both parameters and the return type follow this convention --- \texttt{type [shape] name}. The compiler uses these shapes to verify that every indexing operation is in-bounds at compile time.

Croqtile supports these element types: \texttt{s32} (signed 32-bit int), \texttt{f16} (half-precision float), \texttt{bf16} (brain float), \texttt{f32} (single-precision float), and \texttt{f8\_e4m3} / \texttt{f8\_e5m2} (8-bit floats). In the examples that follow, \texttt{s32} is the simplest to reason about.

\begin{lstlisting}[language=croqtile,caption={Allocating the output}]
s32 [lhs.span] output;
\end{lstlisting}

This allocates the output tensor. The expression \texttt{lhs.span} copies the full shape from \texttt{lhs}, so \texttt{output} automatically has shape \texttt{[4, 8]}. If you later change the input shape, the output adjusts --- a pattern that makes Croqtile functions easier to generalize.

\begin{lstlisting}[language=croqtile,caption={The parallel body}]
parallel {i, j} by [4, 8]
  output.at(i, j) = lhs.at(i, j) + rhs.at(i, j);
\end{lstlisting}

\texttt{parallel \{i, j\} by [4, 8]} creates 4 \(\times\) 8 = 32 parallel instances, one per element position. Each instance runs the body with its own \texttt{(i, j)} pair. The \texttt{.at(i, j)} accessor indexes into a tensor at position \texttt{(i, j)}, and the addition is direct --- one element, one thread.

What \texttt{parallel} leaves unsaid is how those instances map onto the machine: without an annotation the compiler decides. That mapping is deliberately absent from the language here --- the statement describes independent work, not the resources that run it.

\begin{lstlisting}[language=croqtile,caption={Returning the result}]
return output;
\end{lstlisting}

A \texttt{\_\_co\_\_} function returns its result tensor. The return type in the signature (\texttt{s32 [4, 8]}) must match what you actually return. The caller on the host side receives a \texttt{croq:\allowbreak{}:\allowbreak{}spanned\_\allowbreak{}data} --- an owning buffer with shape metadata that you can index into with \texttt{[i][j]}.
